% =====================================================================
% FORMATO DEL PROYECTO INTEGRADOR -- ESTADÍSTICA DESCRIPTIVA (UNL)
% Documento acumulativo: se completa por partes en las tres unidades.
%   Unidad 1: Título, 1. Introducción (problema, pregunta, hipótesis),
%             2. Objetivos y Referencias.
%   Unidad 2: + 3. Metodología y 4. Resultados descriptivos.
%   Unidad 3: + 5. Resultados inferenciales, 6. Discusión y 7. Conclusiones.
%
% Compilación (en este orden):
%   pdflatex proyecto_integrador_estadistica
%   biber    proyecto_integrador_estadistica
%   pdflatex proyecto_integrador_estadistica
%   pdflatex proyecto_integrador_estadistica
% (o, en Overleaf / latexmk, la compilación es automática).
%
% Reglas de la Unidad 1:
%   - La Introducción (secciones 1.1 a 1.3) ocupa como MÁXIMO 2 páginas.
%   - Toda afirmación del planteamiento del problema lleva cita APA:
%       \textcite{clave}  -> Funk et al. (2015)
%       \parencite{clave} -> (Funk et al., 2015)
%   - Las referencias se agregan en referencias.bib y aparecen solas
%     en la sección Referencias cuando se citan en el texto.
%   - Borre todos los textos de instrucción (en gris) antes de entregar.
% =====================================================================
\documentclass[12pt,a4paper]{article}
\usepackage[utf8]{inputenc}
\usepackage[T1]{fontenc}
\usepackage[spanish,es-nodecimaldot]{babel}
\usepackage{newtxtext,newtxmath}            % Times New Roman (APA 7)
\usepackage[margin=2.54cm]{geometry}        % Márgenes de 1 pulgada (APA 7)
\usepackage{setspace}
\usepackage{csquotes}
\usepackage[style=apa,backend=biber]{biblatex}
\addbibresource{referencias.bib}
\usepackage{booktabs,tabularx,array}
\usepackage{graphicx}
\usepackage{caption}
\usepackage{xcolor}
\usepackage{enumitem}
\usepackage{fancyhdr}
\usepackage[hidelinks]{hyperref}

% Encabezado: número de página arriba a la derecha (APA 7)
\pagestyle{fancy}\fancyhf{}
\fancyhead[R]{\thepage}
\renewcommand{\headrulewidth}{0pt}

% Títulos de sección en estilo APA (centrado y negrita / alineado y negrita)
\usepackage{titlesec}
\titleformat{\section}{\centering\normalsize\bfseries}{\thesection.}{0.5em}{}
\titleformat{\subsection}{\normalsize\bfseries}{\thesubsection}{0.5em}{}
\titlespacing*{\section}{0pt}{12pt}{6pt}
\titlespacing*{\subsection}{0pt}{10pt}{4pt}

\setlength{\parindent}{1.27cm}              % Sangría de 0,5 pulgadas (APA 7)
\newcommand{\instruccion}[1]{\textcolor{gray}{\textit{#1}}}

% ---------------- Datos del grupo (complete) ----------------
\newcommand{\titulo}{Título del proyecto: variable o fenómeno, cantón Loja, 2019--2024 (máximo 15 palabras)}
\newcommand{\grupo}{Grupo N.\textsuperscript{o} \_\_ -- Tema del grupo}
\newcommand{\integrantes}{Nombre Apellido \\ Nombre Apellido \\ Nombre Apellido}

\begin{document}

% =====================================================================
% PORTADA (APA 7, versión estudiantil)
% =====================================================================
\begin{titlepage}
\centering
\doublespacing
\vspace*{3cm}
{\bfseries\titulo\par}
\vspace{1.5cm}
\grupo\\[0.4cm]
\integrantes\\[1.2cm]
Carrera de Ingeniería Ambiental, Universidad Nacional de Loja\\
Estadística Descriptiva -- Proyecto integrador\\
Docente: Carlos Guillermo Chuncho Morocho\\
Octubre de 2026
\end{titlepage}

\doublespacing

% =====================================================================
% 1. INTRODUCCIÓN  (máximo 2 páginas: secciones 1.1 a 1.3)
% =====================================================================
\section{Introducción}

\subsection{Planteamiento del problema}

\instruccion{Escriba de 4 a 6 párrafos (alrededor de 1,5 páginas) siguiendo tres pasos. Cada afirmación debe estar respaldada por una cita en normas APA 7. Las oraciones de ejemplo muestran cómo citar; reemplácelas por el contenido de su grupo.}

\instruccion{Paso 1 -- Situación: qué ocurre con el tema de su grupo en el territorio y por qué importa ambientalmente.}
Por ejemplo, la precipitación mensual puede estimarse en zonas con pocas estaciones meteorológicas combinando información satelital y registros de estaciones, como lo hace el conjunto CHIRPS \parencite{funk2015}. De forma similar, la misión Sentinel-2 entrega imágenes multiespectrales de alta resolución que permiten seguir el estado de la vegetación \parencite{drusch2012}, y el producto de detección activa de fuego VIIRS identifica incendios con una resolución de 375~m \parencite{schroeder2014}.

\instruccion{Paso 2 -- Vacío: qué no se conoce, no se ha cuantificado o no se ha descrito para el cantón Loja en el periodo 2019--2024.}
Como señalan \textcite{hernandez2018}, el planteamiento del problema debe delimitar con precisión qué se estudiará, dónde y en qué periodo, de modo que la pregunta pueda responderse con los datos disponibles.

\instruccion{Paso 3 -- Aporte: cómo el análisis estadístico de la base del grupo, obtenida del proyecto FIRELAB-Loja \parencite{firelab2025}, ayuda a llenar ese vacío. El último párrafo conduce directamente a la pregunta de investigación.}

\subsection{Pregunta de investigación}

\instruccion{Una sola pregunta, clara y respondible con la base del grupo. Indique qué variables, dónde (cantón Loja; puede comparar las zonas norte y sur) y cuándo (2019--2024). Evite preguntas de «sí» o «no» y preguntas que pidan demostrar causas.}

¿\ldots{} en el cantón Loja durante el periodo 2019--2024?

\subsection{Hipótesis}

\instruccion{Una afirmación contrastable que anticipa la respuesta. Debe incluir las variables, el lugar, el periodo y un criterio cuantitativo (una proporción, una diferencia entre grupos o una asociación). Se contrastará en la Unidad 3 con inferencia estadística.}

En el cantón Loja, durante el periodo 2019--2024, \ldots{}

% =====================================================================
% 2. OBJETIVOS
% =====================================================================
\section{Objetivos}

\subsection{Objetivo general}

\instruccion{Uno solo. Inicia con un verbo en infinitivo (determinar, caracterizar, analizar, evaluar) y responde directamente la pregunta de investigación.}

\subsection{Objetivos específicos}

\instruccion{Tres objetivos medibles, uno por unidad del curso. Cada uno debe producir un resultado verificable (tabla, gráfico, medida, intervalo o prueba).}

\begin{enumerate}[label=OE\arabic*., leftmargin=1.6cm]
  \item \instruccion{(Unidad 1) Describir y clasificar las variables de \ldots{} mediante tablas de frecuencias.}
  \item \instruccion{(Unidad 2) Caracterizar \ldots{} mediante gráficos, medidas de tendencia central, dispersión y forma, y series temporales.}
  \item \instruccion{(Unidad 3) Estimar / contrastar \ldots{} mediante intervalos de confianza, pruebas chi-cuadrado o regresión logística.}
\end{enumerate}

% =====================================================================
% SECCIONES DE LAS UNIDADES 2 Y 3 (descomente cuando corresponda)
% =====================================================================
% \section{Metodología}                  % Unidad 2
% \subsection{Área de estudio y datos}
% \subsection{Población, muestra y unidad de análisis}
% \subsection{Variables}
% \subsection{Depuración y control de calidad}
% \subsection{Técnicas de análisis}
%
% \section{Resultados descriptivos}      % Unidad 2
%
% \section{Resultados inferenciales}     % Unidad 3
% \section{Discusión}                    % Unidad 3
% \section{Conclusiones}                 % Unidad 3

% =====================================================================
% REFERENCIAS (APA 7, se generan desde referencias.bib)
% =====================================================================
\newpage
\printbibliography[title={Referencias}]

\end{document}
